\chapter{Background \& State of the Art}
\label{sec2: background}
In this section, we provide the biological, computational and mathematical background and context for the 
methods developed in this report. 
We will begin by introducing bacteriophages and the host-prediction problem in Section \ref{sec2.1: host prediction 
problem}, which motivates the need for computational approaches to find a solution. We will then go over the already 
existing and state-of-the-art computational techniques in Section \ref{sec2.2: current approaches}, progressing from the 
classical sequence similarity techniques to the much more efficient and modern protein embeddings techniques. We will 
also look at the limitations that all the currently existing methods have in common, for example the absence of rigorous 
distribution-free uncertainty quantification. In the remaining part of the section, \ref{sec2.3: conformal prediction} 
and \ref{sec2.4: CSA Method} we will introduce the mathematical framework of conformal prediction and the score-based 
aggregation method in \cite{ochoarivera2024conformal}, on which we will build our work.

% However, identifying specific host-range interactions brings some challenges. Traditional laboratory screening via 
% plaque assays is accurate but expensive, labor-intensive, and incapable of scaling to the astronomical number of novel 
% phage genomes generated by modern sequencing \cite{coclet2021global}. Fast, reliable computational prediction tools are 
% therefore essential to prioritize candidate phages prior to experimental validation or to infer host information directly 
% from genomic data.

\section{Bacteriophages and Host-Prediction Problem}
\label{sec2.1: host prediction problem}
Bacteriophages, commonly referred to as phages, are viruses that specifically infect and replicate within 
bacteria. As discussed in Chapter~\ref{sec1:introduction}, the growing threat of antimicrobial resistance and the development 
of phage therapy as an alternative to traditional antibiotics \cite{dejonge2019molecular} have renewed interest in the study 
of phages; beyond therapy, they also play a very important role in shaping microbial communities and evolution in the biosphere.

One of the key characteristics of phages is their host range, which is the set of bacterial hosts that a phage can infect. 
This host range is not fixed for a phage, it is determined by molecular interactions between the phage and the bacterial 
surface structures during the cycle of infection \cite{dejonge2019molecular}, \cite{holtappels2023drivers}. Each phage 
attaches itself to the bacterial host, injects its viral genetic material which are able to replicate inside the host. 
Therefore, some of the phages have a small host range as they can only infect a limited number of 
bacterial strains while there are other which are capable of infecting a broad range of bacterias. 


Identifying which phage infects which host, can have several applications in both environmental 
and clinical studies. In clinical, the host range helps us determine how suitable a phage is for 
designing the treatment for phage therapy. In environmental studies it is valuable to know the 
host being infected to understand the ecological interactions between viruses and microbial 
communities and evolution in the biosphere \cite{holtappels2023drivers}. Precise host range 
knowledge is equally important for avoiding harm, an imprecise or overly broad host assignment 
risks eliminating bacteria that are not the intended target. This may include beneficial species, so 
correctly identifying a phage's host range is as much about protecting non target bacteria as 
it is about targeting the pathogens.

However, host range also brings its own obstacles. In phage therapy, we need to find the exact phage that infects 
the bacteria which is the cause of the infection, and for that we need a fast, reliable host prediction 
method. As noted in Chapter~\ref{sec1:introduction}, traditional laboratory screening is accurate but expensive and time 
consuming, so realistically only a small portion of phage-host pairs can ever be tested, especially since many phages are 
difficult or impossible to culture in the laboratory. Metagenomic sequencing has bypassed this culturing barrier and revealed 
a vast number of new phages directly from the environment; as a result, the number of possible phage-host interactions has 
grown astronomically and computational predictions have become essential \cite{coclet2021global}.

% Compounding this, isolating and culturing a phage in the first place is itself far from trivial, since many phages 
% depend on hosts that are difficult or impossible to culture under standard laboratory conditions, which limits how 
% many phage-host pairs can even be brought forward for testing. The advent of metagenomic sequencing has partly bypassed 
% this bottleneck, uncovering vast numbers of previously uncharacterized phage genomes directly from environmental samples 
% without requiring culturing or isolation at all. Given this rapid, culture-independent increase in the number of known 
% phage genomes, the number of possible phage-host interactions has grown astronomically, and this leads us to the need 
% for a computational prediction method \cite{coclet2021global}.

We formulate the host prediction problem as follows: let $\mathcal{X} \subseteq \mathbb{R}^{d}$ denote the feature space 
of phage representations, and let $\mathcal{H} = \{h_{1}, \dots, h_{N}\}$ be the discrete label space of candidate bacterial
hosts.

In a single-host classification setting, we observe random pairs $(X, H) \sim \mathbb{P}_{X, H}$, where $X \in \mathcal{X}$ is 
a feature vector representing a phage and $H \in \mathcal{H}$ denotes its true host. In a generalized multi-host setting, the 
target is a non-empty subset $H^* \subseteq \mathcal{H}$, i.e., $H^* \in \mathcal{P}(\mathcal{H}) \setminus \{\emptyset\}$, 
where $\mathcal{P}(\mathcal{H}) = 2^{\mathcal{H}}$ denotes the power set of $\mathcal{H}$.

Our primary objective is to construct a set valued prediction mapping
\begin{equation}
C: \mathcal{X} \longrightarrow \mathcal{P}(\mathcal{H})
\end{equation}
which maps an unlabelled test phage $X \in \mathcal{X}$ to a prediction set $C(X) \subseteq \mathcal{H}$ satisfying the finite 
sample coverage guarantee at a user specified significance level $\alpha \in (0, 1)$:
\begin{equation}
\mathbb{P}(H \in C(X)) \ge 1 - \alpha.
\end{equation}

where $\alpha \in (0,1)$ denotes the target error probability.
